\section*{2 Subset-Wise Symmetric Polynomials}
In this section, we will formally define subset-wise symmetric polynomials and prove suitable tail bounds for them. This section can be read on its own, independent of the application to PRGs. We start by recalling known tail bounds for elementary symmetric polynomials.
\subsection*{2.1 Gopalan and Yehudayoff's Bounds for Symmetric Polynomials}
As a reminder, the $k$-th elementary symmetric polynomial is defined by
\[
S_k(y)=\sum_{\substack{I\subseteq[m]\\|I|=k}}\prod_{i\in I}y_i.
\]
We rely on the following tail bound by Gopalan and Yehudayoff [17]. As discussed in Section1.5.1, the bound says that if two $S_k$-s in a row are small, then all subsequent $S_k$-s are small.
\begin{authorclaim}{Theorem2 ([17])}
Let $y\in\mathbb R^m$, $\theta>0$, and $\ell\in\mathbb N$ satisfy $S_\ell^2(y)\le\theta^\ell/\ell!$ and $S_{\ell+1}^2(y)\le\theta^{\ell+1}/(\ell+1)!$. Then, for every $k\ge\ell$,
\[
|S_k(y)|\le\left(\frac{64e^2\theta\ell}{k}\right)^{k/2}.
\]
\end{authorclaim}
The exact statement of Theorem2 does not appear in Gopalan and Yehudayoff's work [17], but it follows readily from their analysis, and it was used previously by Meka et al. [24, Theorem5.2].\footnote{The careful reader will notice a slight discrepancy between the exact constants of Theorem2 on the one hand and the statements by Gopalan and Yehudayoff [17] and Meka et al. [24] on the other. This discrepancy reflects a minor mistake in the original paper by Gopalan and Yehudayoff [17] that we have here corrected.}
\subsection*{2.2 Our Tail Bounds for Subset-Wise Symmetric Polynomials}
Let $B=(B_1,\ldots,B_Q)$ be a partition of $[m]$, namely $[m]=B_1\sqcup\cdots\sqcup B_Q$. (The sets $B_1,\ldots,B_Q$ correspond to the "buckets" discussed in Section1.5.1.) Throughout this paper, let $\mathbb N$ denote the set of nonnegative integers, $\mathbb N=\{0,1,2,\ldots\}$. For a vector $\vec k=(\vec k[1],\ldots,\vec k[Q])\in\mathbb N^Q$ and $y\in\mathbb R^m$, we define the following polynomial:
\[
S_{\vec k,B}(y)=\sum_{\substack{I\subseteq[m]\\\forall q,\ |B_q\cap I|=\vec k[q]}}\prod_{i\in I}y_i.
\]
We name these polynomials as subset-wise symmetric polynomials, since for every $q\in[Q]$, $S_{\vec k}(y)$ when restricted to the $B_q$ variables is a degree $\vec k[q]$ symmetric polynomial.
Throughout this section we fix $B=(B_1,\ldots,B_Q)$ to be a partition of $[m]$. When the partition $B$ is clear from the context, we will simply write $S_{\vec k}$ instead of $S_{\vec k,B}$. To formulate our tails bounds for the subset-wise symmetric polynomials, we will need the following auxiliary polynomials:
\[
R_{\vec k}(y)\mathrel{\mathop:}=S_{\vec k}^2(y)\cdot\prod_{q=1}^Q\vec k[q]!.
\]
Given $c>1$, we will assign each vector $\vec k\in\mathbb N^Q$ a weight, defined as
\[
\|\vec k\|_{(c)}=\sum_{q=1}^Qc^q\vec k[q].
\]
(In our PRG application, $B_q$ will be the set of terms with approximately $c^q$ input variables, so $\|\vec k\|_{(c)}$ will be approximately the number of input variables participating in each monomial of $S_{\vec k}$, as outlined in Section1.5.1.) It is easy to verify that the above weight function is indeed a norm; however, we will not be using this observation.
The main result of this section is a tail-bound for subset-wise symmetric polynomials. In Lemma3, the parameter $A$ is analogous to the "cutoff point" $k_0$ discussed in Section1.5.1.
\begin{authorclaim}{Lemma3}
Suppose $c>1$ and $Q,A\in\mathbb N$ satisfy $A>\max\{(10^6c/(c-1))\cdot c^Q,2^{60}Q^2\}$. Let $Y$ be a random variable taking values in $\mathbb R^m$. Moreover, suppose for every $\vec k\in\mathbb N^Q$ with $\|\vec k\|_{(c)}\le A$,
\[
\mathbb E_Y[R_{\vec k}(Y)]\le2^{-\|\vec k\|_{(c)}/8}.
\]
Then, except with probability $2^{-A/2^{23}}$ over $y\sim Y$,
\[
\sum_{\substack{\vec k\in\mathbb N^Q\\\|\vec k\|_{(c)}>A}}|S_{\vec k}(y)|\le2^{-A/1024}.
\]
\end{authorclaim}
Lemma3 is similar in spirit to Theorem2: it says that if the "early" subset-wise symmetric polynomials are small (with high probability), then the "late" subset-wise symmetric polynomials are all small (with high probability).
\subsection*{2.3 Non-probabilistic Tail Bound}
Before moving to the proof of Lemma3 in the next subsection, here we first give a tail-bound in the case when the input $y$ satisfies some useful properties. We will later prove Lemma3, by showing that a random $Y$ satisfies these properties with high probability. Given a vector $\vec k\in\mathbb N^Q$, we define the restriction of $\vec k$ to a set $\mathcal Q\subseteq[Q]$ by
\[
\vec k|_{\mathcal Q}[q]=\begin{cases}\vec k[q]&q\in\mathcal Q,\\0&q\notin\mathcal Q.\end{cases}
\]
Our non-probabilistic tail bound goes as follows.
\begin{authorclaim}{Lemma4}
Suppose $c>1$ and $Q,A\in\mathbb N$ satisfy $A>\max\{(10^6c/(c-1))\cdot c^Q,2^{60}Q^2\}$. Let $y\in\mathbb R^m$ be a fixed vector. Suppose that for every $\vec k\in\mathbb N^Q$, with $A/10^5\le\|\vec k\|_{(c)}\le A$, and for every pair of disjoint sets $\mathcal Q_1,\mathcal Q_2\subseteq[Q]$ satisfying $\{q:\vec k[q]>1\}\subseteq\mathcal Q_1\cup\mathcal Q_2$, we have
\[
R_{\vec k|_{\mathcal Q_1}}(y)\cdot R_{\vec k|_{\mathcal Q_2}}(y)^4\le2^{-\|\vec k\|_{(c)}/32}.
\]
Then,
\[
\sum_{\substack{\vec k\in\mathbb N^Q\\\|\vec k\|_{(c)}>A}}|S_{\vec k}(y)|\le2^{-A/1024}.
\]
\end{authorclaim}
\begin{proof}
For a fixed $\ell\in\mathbb N$ and $q\in[Q]$, define
\[
S_{\ell,q}=\sum_{\substack{I\subseteq B_q\\|I|=\ell}}\prod_{i\in I}y_i,
\]
which is the $\ell$-th elementary symmetric polynomial applied to $(y_i)_{i\in B_q}$. Similarly, define $R_{\ell,q}=S_{\ell,q}^2(y)\cdot\ell!$.
Fix $\vec k$ with $\|\vec k\|_{(c)}>A$, let $\lambda=10^5\cdot\|\vec k\|_{(c)}/A$ and let $\vec k'\in\mathbb N^Q$ be such that $\vec k'[q]=\lceil\vec k[q]/\lambda\rceil$. Thus, $A/10^5\le\|\vec k'\|_{(c)}\le A/2$. Let $\mathcal Q:=\{q\in Q:\vec k[q]\ge1\}$, and for each $q\in\mathcal Q$, let $\theta_q>0$ be the smallest\footnote{It is possible that $\theta_q=0$ satisfies Equation(3). In this degenerate case, we must have $S_{\vec k[q],q}=0$. This implies $S_{\vec k}(y)=0$, hence Equation(4) trivially holds.} number satisfying
\begin{equation}
R_{\vec k'[q],q}\le\theta_q^{\vec k'[q]}
\quad\text{and}\quad
R_{\vec k'[q]+1,q}\le\theta_q^{\vec k'[q]+1}.\tag{3}
\end{equation}
By Theorem2,
\[
|S_{\vec k[q],q}|\le\left(\frac{64e^2\theta_q\vec k'[q]}{\vec k[q]}\right)^{\vec k[q]/2}.
\]
Subset-wise symmetric polynomials by design can be expressed as a product of elementary symmetric polynomials, hence
\begin{align*}
|S_{\vec k}(y)|&=\prod_{q=1}^Q|S_{\vec k[q],q}(y)|
\le\prod_{q\in\mathcal Q}\left(\frac{64e^2\theta_q\vec k'[q]}{\vec k[q]}\right)^{\vec k[q]/2}\\
&=\left(\prod_{q\in\mathcal Q}\theta_q^{\vec k[q]/\lambda}\right)^{\lambda/2}
\cdot\prod_{q\in\mathcal Q}\left(8e\sqrt{\vec k'[q]/\vec k[q]}\right)^{\vec k[q]}.
\end{align*}
By our choice of $\theta_q$,
\begin{align*}
\theta_q^{\vec k'[q]}
&=\max\{R_{\vec k'[q],q}(y),R_{\vec k'[q]+1,q}(y)^{\vec k'[q]/(\vec k'[q]+1)}\}\\
&\le\max\{R_{\vec k'[q],q}(y),R_{\vec k'[q]+1,q}(y),\sqrt{R_{\vec k'[q]+1,q}(y)}\}.
\end{align*}
Observe that $\vec k[q]/\lambda\in[\vec k'[q]-1,\vec k'[q]]$, and thus $\theta_q^{\vec k[q]/\lambda}$ is between $\theta_q^{\vec k'[q]-1}$ and $\theta_q^{\vec k'[q]}$. If $\vec k'[q]=1$, then $\theta_q^{k'_q-1}=1$, and otherwise $\theta_q^{\vec k'[q]-1}$ is between $\theta_q^{\vec k'[q]}$ and $\sqrt{\theta_q^{\vec k'[q]}}$. Therefore,
\begin{align*}
\theta_q^{\vec k[q]/\lambda}
&\le\max\{\theta_q^{\vec k'[q]},\sqrt{\theta_q^{\vec k'[q]}},1_{\vec k'[q]=1}\}\\
&\le\max\{R_{\vec k'[q],q}(y),R_{\vec k'[q],q}(y)^{1/4},R_{\vec k'[q]+1,q}(y),R_{\vec k'[q]+1,q}(y)^{1/4},1_{\vec k'[q]=1}\}.
\end{align*}
For every $q$, choose $\vec k''[q]\in\{\vec k'[q],\vec k'[q]+1\}$ such that
\[
\theta_q^{\vec k[q]/\lambda}\le\max\{R_{\vec k''[q],q}(y),R_{\vec k''[q],q}(y)^{1/4},1_{\vec k''[q]=1}\}.
\]
Note that $\|\vec k''\|_{(c)}\ge\|\vec k'\|_{(c)}$ and
\[
\|\vec k''\|_{(c)}\le\|\vec k'\|_{(c)}+\sum_{q=1}^Qc^q\le\|\vec k'\|_{(c)}+\frac A{10^6}<A.
\]
Therefore, there exist disjoint sets $\mathcal Q_1,\mathcal Q_2\subseteq[Q]$ such that $\{q:\vec k''[q]>1\}\subseteq\mathcal Q_1\cup\mathcal Q_2$, and that for every $q\in\mathcal Q$,
\[
\theta_q^{\vec k[q]/\lambda}\le
\begin{cases}R_{\vec k''[q],q}(y)&q\in\mathcal Q_1,\\R_{\vec k''[q],q}(y)^{1/4}&q\in\mathcal Q_2,\\1&\text{otherwise}.\end{cases}
\]
Multiplying over $q\in\mathcal Q$, we get
\begin{align*}
\prod_{q\in\mathcal Q}\theta_q^{\vec k[q]/\lambda}
&\le\prod_{q\in\mathcal Q_1}R_{\vec k''[q],q}(y)\cdot\prod_{q\in\mathcal Q_2}R_{\vec k''[q],q}(y)^{1/4}\\
&=\left(R_{\vec k''|_{\mathcal Q_1}}(y)^4\cdot R_{\vec k''|_{\mathcal Q_2}}(y)\right)^{1/4}
\le2^{-\|\vec k''\|_{(c)}/128}\le2^{-\|\vec k'\|_{(c)}/128}.
\end{align*}
As a result,
\begin{align}
|S_{\vec k}(y)|
&\le2^{-\frac{\|\vec k'\|_{(c)}}{128}\cdot\frac\lambda2}
\cdot\prod_{q\in\mathcal Q}\left(8e\sqrt{\vec k'[q]/\vec k[q]}\right)^{\vec k[q]}\notag\\
&\le2^{-\frac{\|\vec k'\|_{(c)}}{256}\lambda}\cdot
\prod_{q\in\mathcal Q}\left(8e\sqrt{2/10^5}\right)^{\vec k[q]}\cdot\left(\sqrt{10^5/2}\right)^\lambda\notag\\
&\le2^{-\frac{\|\vec k'\|_{(c)}}{256}\lambda}\cdot2^{8Q\lambda}\cdot4^{-\|\vec k\|_1}
\le2^{-\frac{\|\vec k'\|_{(c)}}{512}\lambda}\cdot4^{-\|\vec k\|_1}
\le2^{-\|\vec k\|_{(c)}/512}\cdot4^{-\|\vec k\|_1}.\tag{4}
\end{align}
To see the second inequality, observe that when $\vec k[q]>\lambda$, then
$\left(8e\sqrt{\vec k'[q]/\vec k[q]}\right)^{\vec k[q]}
\le(8e\sqrt{2/10^5})^{\vec k[q]}$, and otherwise
$\left(8e\sqrt{\vec k'[q]/\vec k[q]}\right)^{\vec k[q]}\le(8e)^\lambda$.
Summing up over all choices of $\vec k$ we get,
\begin{align*}
\sum_{\substack{\vec k\in\mathbb N^Q\\\|\vec k\|_{(c)}>A}}|S_{\vec k}(y)|
&\le\sum_{L=1}^m\sum_{\substack{\vec k\in\mathbb N^Q\\\|\vec k\|_{(c)}>A,\|\vec k\|_1=L}}2^{-\|\vec k\|_{(c)}/512}\cdot4^{-L}\\
&\le2^{-A/512}\cdot\sum_{L=1}^m4^{-L}\cdot|\{\vec k\in\mathbb N^q:\|\vec k\|_1=L\}|\\
&=2^{-A/512}\cdot\sum_{L=1}^m4^{-L}\cdot\binom{Q-1+L}{Q-1}\\
&\le2^{-A/512}\cdot\sum_{L=1}^m4^{-L}\cdot2^{Q-1+L}
\le2^{-A/512}\cdot2^{Q-1}\cdot\sum_{L=1}^m2^{-L}\le2^{-A/1024}.
\end{align*}
\end{proof}
\subsection*{2.4 Probabilistic Tail Bound: Proof of Lemma3}
\begin{proof}
Let $\vec k,\mathcal Q_1$, and $\mathcal Q_2$ be as in the statement of Lemma4. Using the Cauchy-Schwarz inequality and the concavity of $(\cdot)^{1/4}$, we get
\begin{align*}
\mathbb E[(R_{\vec k|_{\mathcal Q_1}}(Y))^{1/8}(R_{\vec k|_{\mathcal Q_2}}(Y))^{1/2}]
&\le\left(\mathbb E[(R_{\vec k|_{\mathcal Q_1}}(Y))^{1/4}]\cdot\mathbb E[R_{\vec k|_{\mathcal Q_2}}(Y)]\right)^{1/2}\\
&\le\left(\mathbb E[R_{\vec k|_{\mathcal Q_1}}(Y)]^{1/4}\cdot\mathbb E[R_{\vec k|_{\mathcal Q_2}}(Y)]\right)^{1/2}\\
&\le(2^{-\|\vec k|_{\mathcal Q_1}\|_{(c)}/32}\cdot2^{-\|\vec k|_{\mathcal Q_2}\|_{(c)}/8})^{1/2}\\
&\le2^{-\|\vec k|_{\mathcal Q_1\cup\mathcal Q_2}\|_{(c)}/64}\\
&\le2^{-(\|\vec k\|_{(c)}-(c/(c-1))c^Q)/64}\\
&\le2^{-(\|\vec k\|_{(c)}-A/20000)/64}\le2^{-\|\vec k\|_{(c)}/128}.
\end{align*}
Therefore, by Markov's inequality, except with probability at most $2^{-\|\vec k\|_{(c)}/256}\le2^{-A/2560000}$, we have
\[
R_{\vec k|_{\mathcal Q_1}}(Y)\cdot R_{\vec k|_{\mathcal Q_2}}(Y)^4\le2^{-\|\vec k\|_{(c)}/32}.
\]
The above analysis was done for a fixed choice of $\vec k,\mathcal Q_1$, and $\mathcal Q_2$. The number of choices for such $\vec k$ is $A^Q$ (which is subexponential in $A$), and the number of such $\mathcal Q_1,\mathcal Q_2$ is at most $3^Q$ (which is a polynomial in $A$), thus Lemma3 follows by a union bound. More precisely, one can check that since $A\ge2^{60}Q^2$, then $(3A)^Q\cdot2^{-A/2560000}\le2^{-A/2^{23}}$.
\end{proof}
